\documentclass[conference]{IEEEtran}

\usepackage{cite}
\usepackage{amsmath,amssymb,amsfonts}
\usepackage{algorithmic}
\usepackage{graphicx}
\usepackage{textcomp}
\usepackage{xcolor}

\begin{document}

\title{R-LAM: Reproducibility-Constrained Large Action Models for Scientific Workflow Automation}

\author{
    \IEEEauthorblockN{Suriya Sureshkumar\textsuperscript{1,*} Ivan Nilash X\textsuperscript{1,*}}
    \IEEEauthorblockA{\textsuperscript{1,*}Department of AI \& Data Science, RMK Engineering College, Chennai, India}
    \IEEEauthorblockA{\{240171.ad, 240227.ad\}@rmkec.ac.in}
    \thanks{*Both authors contributed equally to this work.}
}

\maketitle

\begin{abstract}
Large Action Models (LAMs) extend large language models by enabling autonomous decision-making and tool execution, making them promising for automating scientific workflows. However, scientific workflows impose strict requirements on reproducibility, auditability, and deterministic execution, which are not satisfied by generic LLM-based agents. Unconstrained action generation can lead to silent state changes, non-deterministic executions, and irreproducible experimental results, limiting the applicability of LAMs in scientific settings.

In this paper, we propose R-LAM, a reproducibility-constrained framework for applying Large Action Models to scientific workflow automation. R-LAM introduces structured action schemas, deterministic execution policies, and explicit provenance tracking to ensure that every action and intermediate artifact is auditable and replayable. The framework supports failure-aware execution loops and controlled workflow forking, enabling iterative experimentation without compromising reproducibility.

We implement R-LAM as a lightweight Python framework and release it as an open-source PyPI package to facilitate reproducible research. Experimental evaluation on representative scientific workflows demonstrates that R-LAM improves reproducibility success rates and execution reliability compared to unconstrained LLM-based agents, while retaining adaptive control over workflow execution.
\end{abstract}

\begin{IEEEkeywords}
Reproducibility, Large Action Models, Scientific Workflows, Provenance Tracking
\end{IEEEkeywords}

\section{Introduction}

Large language models (LLMs) have demonstrated strong capabilities in reasoning, code generation, and tool usage, enabling a new class of autonomous systems commonly referred to as agents or Large Action Models (LAMs). Unlike traditional LLMs, LAMs are designed to select and execute actions such as running code, invoking APIs, or orchestrating multi-step workflows. These capabilities have motivated growing interest in applying LAMs to automate scientific data analysis pipelines and experimental workflows.

Despite this promise, scientific workflows impose constraints that fundamentally differ from general-purpose automation. Scientific experiments are expected to be reproducible, auditable, and deterministic, such that results can be independently verified and extended by other researchers. Generic LLM-based agents often operate with implicit state, unconstrained action generation, and stochastic execution behavior, which can introduce silent state changes, non-deterministic outcomes, and incomplete provenance records.

To address these challenges, we introduce R-LAM, a reproducibility-constrained framework for applying Large Action Models to scientific workflow automation. R-LAM treats actions as structured, first-class entities, enforces deterministic execution policies, and records complete execution traces for provenance and replay. The framework enables adaptive experimentation through failure-aware execution and controlled workflow forking, while preserving reproducibility guarantees.

\section{Challenges in Scientific Workflow Reproducibility}

Reproducibility remains a central challenge in computational and data-intensive scientific research. Prior studies show that many published computational results cannot be independently reproduced due to undocumented software environments, missing parameters, and incomplete provenance capture. Beaulieu-Jones \textit{et al.} demonstrate that even when code and data are available, substantial engineering effort is often required to reproduce results, motivating automated reproducibility mechanisms \cite{beaulieu2017reproducibility}.

Scientific workflows frequently contain hidden state and implicit dependencies that are not captured in workflow specifications or execution logs. Environment configurations, external tool versions, and ephemeral intermediate artifacts may influence results without being explicitly recorded. Suetake \textit{et al.} highlight the difficulty of validating workflow reproducibility in bioinformatics, even in containerized settings, due to the lack of automated equivalence verification mechanisms \cite{suetake2023workflow}.

Non-determinism further exacerbates reproducibility challenges. Stochastic algorithms, parallel execution, and hardware variability can cause repeated runs with identical inputs to yield divergent outputs. Horton \textit{et al.} emphasize that simple version control of scripts is insufficient for reproducibility, underscoring the need for comprehensive environment and provenance management \cite{horton2022reproducibility}. In contrast to business automation, scientific workflows must treat every execution as a first-class scientific artifact.

\section{Related Works}

\subsection{LLM Agents for Scientific Workflows}

Recent work explores using LLMs to assist scientific workflow development. Shin \textit{et al.} describe the evolution of scientific workflows in the agentic AI era, arguing that LLM-driven agents can progressively automate scientific processes \cite{shin2025agentic}. Gupta \textit{et al.} demonstrate that ChatGPT can assist users in understanding and authoring scientific workflows, though primarily in a human-in-the-loop setting \cite{gupta2023chatgpt}. FlowMind, proposed by Zeng \textit{et al.}, uses LLMs to generate structured workflows from high-level task descriptions \cite{zeng2024flowmind}. Masera \textit{et al.} introduce Snakemaker, which converts ad-hoc scripts into maintainable Snakemake workflows \cite{masera2025snakemaker}.

Beyond workflow generation, Large Action Models formalize agentic execution. Dodge \textit{et al.} define LAMs as LLM-based systems that generate and execute actions in external environments \cite{dodge2024lam}. Sangchai \textit{et al.} extend this paradigm to human-in-the-loop robotic systems \cite{sangchai2025robotlam}. Benchmarks such as AgentBench \cite{liu2023agentbench}, TheAgentCompany \cite{xu2025agentcompany}, and MedAgentBench \cite{jiang2025medagentbench} show that while LLM agents can perform multi-step reasoning, their execution behavior remains stochastic and difficult to control.

\subsection{Workflow Management Systems and Reproducibility}

Workflow management systems emphasize deterministic execution and provenance tracking. Beaulieu-Jones and Greene propose continuous analysis to automate reproducibility \cite{beaulieu2017reproducibility}. Mione \textit{et al.} present a DAG-based workflow system for reproducible self-driving laboratory experiments \cite{mione2024workflow}. Grayson \textit{et al.} study reproducibility at scale by re-executing Snakemake and nf-core workflows \cite{grayson2023snakemake}. WorkflowHub provides an open registry for sharing computational workflows \cite{goble2025workflowhub}. Domain-specific systems such as Atomate2 \cite{rich2025atomate2} and AlabOS \cite{fei2024alabos} further demonstrate reproducible workflow execution but rely on predefined pipelines without adaptive reasoning.

\subsection{Automated Experiment Orchestration and AutoML}

Automation of experimental design has also been explored. Gierisch and Mauerer introduce QEF, a framework for reproducible quantum software experiments \cite{gierisch2025qef}. Xu \textit{et al.} propose an AutoML-based workflow for design-of-experiments selection and benchmarking \cite{xu2024automl}. These systems automate parameter exploration but operate within fixed execution structures.

\subsection{Gap Analysis}

Existing approaches either emphasize autonomous reasoning without execution guarantees or deterministic workflows without adaptive control. Our work bridges this gap by treating reproducibility as a core constraint within LAM-based execution, enabling adaptive scientific workflows while preserving auditability, determinism, and replayability. Rather than extending a single prior system, our work synthesizes insights from LLM-based action models and reproducible workflow management to address execution-level gaps in scientific automation.

\section{Large Action Models: Opportunities \& Risks}

Large Action Models (LAMs) extend large language models by coupling high-level reasoning with the ability to execute concrete actions in external environments. This execution-centric paradigm introduces new opportunities for automating complex scientific workflows, but also exposes fundamental risks when applied to scientific settings where execution semantics, provenance, and determinism are critical. This section analyzes both aspects to motivate the need for reproducibility-aware execution constraints.

\subsection{Opportunities}

A central advantage of LAMs is their ability to support \emph{adaptive control}. Unlike static workflow engines, LAMs can condition future actions on intermediate results, enabling workflows that respond dynamically to observed outcomes. This allows scientific pipelines to adjust execution paths when unexpected data characteristics, quality issues, or partial failures are encountered.

LAMs also enable \emph{dynamic parameter tuning} during execution. Rather than fixing experimental parameters in advance, an LAM can iteratively refine configurations based on observed performance signals. This capability is particularly useful in exploratory scientific workflows, where optimal parameters are not known a priori and must be discovered through iterative experimentation.

Another key opportunity is \emph{failure-aware replanning}. Traditional workflow systems typically halt execution upon failure or require manual intervention. In contrast, LAMs can reason about failures, infer potential causes, and attempt alternative execution strategies, such as rerunning steps with modified inputs or substituting tools. This adaptivity has the potential to reduce human effort in managing long-running and complex scientific workflows.

\subsection{Risks}

Despite these advantages, the execution semantics of Large Action Models remain poorly defined in scientific settings. Because action selection is driven by probabilistic model outputs and executed through heterogeneous external tools, LAM-based systems lack inherent guarantees about state isolation, execution ordering, and outcome consistency.

One major risk is \emph{action hallucination}, where the model generates invalid, unsafe, or non-existent actions. Unlike hallucinated text, hallucinated actions can directly modify external state, corrupt datasets, or invalidate experimental results, often without immediate detection.

LAM-driven execution is also susceptible to \emph{non-deterministic behavior}. Variability in model outputs, stochastic decoding processes, and interactions with external systems can cause repeated executions with identical inputs to diverge. Such variability undermines the repeatability required for scientific validation and comparison.

Another critical concern is the presence of \emph{undocumented state changes}. LAMs may implicitly depend on or modify hidden state, such as temporary files, environment variables, or external service configurations, without explicitly recording these changes. These hidden dependencies complicate auditing, replay, and independent verification of results.

Collectively, these factors lead to \emph{result irreproducibility}, where outcomes cannot be reliably regenerated even when code and data are preserved.

\subsection{Implications for Scientific Automation}

Without explicit constraints, LAMs optimize for task completion rather than scientific validity. While this objective may be acceptable in general automation or business process settings, it is fundamentally misaligned with the requirements of scientific research, which demand determinism, auditability, and complete provenance. Addressing these risks is therefore a prerequisite for responsibly deploying LAMs in scientific workflow automation.

\section{Reproducibility-Constrained LAM Framework}

\begin{figure}[ht!]
    \centering
    \includegraphics[width=0.5\textwidth]{R-LAM_HL_Architecture.png}
    \caption{High-level architecture of R-LAM. The Large Action Model proposes structured actions that are validated and executed by a deterministic engine. All execution effects are recorded in a provenance-aware trace store, which provides observable state for subsequent reasoning.}
    \label{R-LAM_HL_Architecture}
\end{figure}

This section presents the core technical contribution of this work: a reproducibility-constrained execution framework for Large Action Models (LAMs). The framework is designed to enable adaptive, agent-driven workflow execution while enforcing strict guarantees on auditability, determinism, and replayability. Rather than modifying the internal reasoning of the LAM, the framework constrains execution at the action level, treating reproducibility as a first-class system invariant.

\subsection{Action Schema}

At the foundation of the framework is a formal definition of an \emph{action}. An action represents the smallest unit of executable behavior issued by a LAM and is defined as a structured, immutable object:
\[
\begin{aligned}
a = \langle\;& id,\ type,\ inputs,\ parameters,\\
             & preconditions,\ effects,\ metadata \;\rangle
\end{aligned}
\]
where $id$ uniquely identifies the action instance, $type$ specifies the execution primitive (e.g., tool invocation, script execution), $inputs$ and $parameters$ define all required data and configuration values, $preconditions$ encode required state constraints, $effects$ describe expected state transitions, and $metadata$ captures contextual information such as timestamps, environment identifiers, and model configuration.

Actions are strictly declarative: they describe \emph{what} is to be executed, not \emph{how} execution is performed. This separation ensures that all executable intent is explicitly represented before execution, eliminating implicit or hidden behaviors.

\subsection{Execution Engine}

The execution engine serves as a deterministic mediator between the LAM and the external environment. Rather than executing actions directly, the engine validates each proposed action against its schema, checks preconditions, and enforces execution policies before dispatch.

To ensure determinism, the engine enforces a fixed execution order, controlled randomness policies, and explicit environment binding. External side effects are isolated through sandboxed execution contexts, and all tool invocations are mediated by adapters that normalize inputs and outputs.

By decoupling action selection from action execution, the framework allows the LAM to remain adaptive while preventing uncontrolled state changes. The execution engine thus acts as a reproducibility firewall between probabilistic reasoning and deterministic execution.

We assume that external tools invoked by the execution engine behave deterministically with respect to recorded inputs and environment bindings.

\subsection{Provenance and Trace Logging}

\begin{figure*}[t]
    \centering
    \includegraphics[width=1\textwidth]{Execution_Trace_Graph_(DAG).png}
    \caption{Illustrates an execution trace graph generated by the framework. Each node corresponds to a single executed action, including its inputs, outputs, execution environment binding, and terminal status. Nodes marked with a failure indicator represent actions that terminated unsuccessfully but were nevertheless fully logged and retained in the trace.
    Importantly, failed actions are not removed, overwritten, or retried implicitly. Instead, failures become first-class nodes in the execution graph, enabling subsequent recovery actions to be explicitly linked through control dependencies. This design ensures that execution history remains complete and auditable, and that all downstream results can be traced to a concrete sequence of observed events.
}
    \label{fig:Execution_Trace_Graph_(DAG)}
\end{figure*}

\begin{figure*}[t]
    \centering
    \includegraphics[width=1\textwidth]{Replay_and_Forking_Mechanism.png}
    \caption{Illustrates the replay and forking mechanism in R-LAM. A new execution trace is derived from an existing trace by reusing a shared prefix of logged actions without re-execution. Prior actions are replayed by reusing their persisted outputs, ensuring identical inputs across experiments. Forking introduces new actions only at the divergence point, where modified parameters are explicitly recorded, while the original execution trace remains immutable. This design enables controlled exploration of alternative hypotheses while preserving auditability, determinism, and complete provenance.
    Unlike naive reruns, replay prevents nondeterminism arising from repeated execution under evolving environments.
}
    \label{Replay_and_Forking_Mechanism}
\end{figure*}

Every executed action produces a provenance record that is appended to a global execution trace. The trace is represented as a directed acyclic graph (DAG), where nodes correspond to action instances, and edges encode data and control dependencies between actions.

Each node stores the complete action schema, execution status, outputs, environment hashes, and timestamps. Dependencies explicitly capture how artifacts flow between actions, enabling downstream analysis of execution lineage.

\textbf{In our framework, an action that is not logged is treated as non-existent.}  
This invariant ensures that all state transitions contributing to a result are observable, auditable, and reproducible. Any side effect not reflected in the execution trace is considered invalid by construction.

\subsection{Replay and Forking}

The execution trace graph enables two critical operations: \emph{replay} and \emph{forking}. Replay reconstructs a prior workflow execution by reconstructing execution outcomes by reusing logged action outputs in topological order using the logged inputs, parameters, and environment bindings. Because all actions are deterministic with respect to their recorded context, replay yields identical results under equivalent execution conditions.

Forking allows controlled divergence from a prior execution. Given a trace prefix, the system can spawn a new execution branch by modifying selected action parameters or inputs while preserving all upstream provenance. This enables exploratory experimentation without contaminating the original execution lineage, supporting reproducible what-if analysis.

Both replay and forking operate solely on the execution trace, requiring no access to the original LAM reasoning process.

\subsection{Failure Handling}

Failures are treated as first-class events within the framework. When an action fails, the execution engine records the failure type, error context, and partial outputs in the trace. Rather than terminating execution unconditionally, the framework enters a failure-aware control loop.

Within this loop, the LAM may propose alternative actions, parameter adjustments, or execution paths. Each recovery attempt is logged as a new action node, preserving a complete record of failure and recovery behavior. This design enables post hoc analysis of failure modes and prevents silent recovery that would otherwise compromise reproducibility.

By explicitly modeling failure and recovery within the execution trace, the framework supports adaptive behavior while maintaining strict guarantees on auditability and replayability.

\section{Implementation \& PyPI Artifact}

To support reproducibility and facilitate independent evaluation, we implement the proposed framework as a lightweight, open-source software artifact. The implementation is designed to closely reflect the conceptual framework described in Section~V, prioritizing determinism, auditability, and traceability over performance optimization or feature completeness.

\subsection{Implementation Details}

The framework is implemented in Python, chosen for its widespread adoption in scientific computing and compatibility with existing data analysis ecosystems. The system is organized into modular components corresponding to the core framework abstractions, including action schemas, the execution engine, and provenance management. Clear interfaces separate LAM reasoning from execution, enabling the framework to operate with different LLM backends without modification to the execution layer.

Deterministic execution is enforced through explicit control of execution order, environment binding, and randomness sources. All external tool invocations are mediated by adapter layers that normalize inputs and outputs and restrict side effects. Where stochastic components are unavoidable, configuration parameters and random seeds are explicitly recorded as part of the action metadata.

\subsection{Packaging and Distribution}

The framework is packaged as a standard Python package and distributed via the Python Package Index (PyPI). The package includes a minimal command-line interface and a programmatic API that allows users to define actions, execute workflows, and inspect execution traces. Dependencies are explicitly specified to reduce environment-related variability, and the package can be installed using standard Python tooling.

For peer review, the artifact can be released in an anonymized repository, following common reproducibility practices. The final version of the package will be made publicly available upon acceptance, enabling reuse and extension by the research community.

\subsection{Intended Scope}

The implementation is intended as a lightweight research artifact rather than a full-featured workflow engine. It is not designed to replace existing workflow management systems such as Airflow or Snakemake, but to demonstrate how reproducibility constraints can be systematically integrated into LAM-based execution. The focus is on clarity, correctness, and traceability, providing a concrete reference implementation that complements the conceptual framework presented in this paper.

\section{Limitations \& Ethical Considerations}

This work is intended as a systems-level exploration of reproducibility-constrained execution for Large Action Models in scientific workflows. While the proposed framework demonstrates the feasibility and benefits of enforcing reproducibility constraints at the action level, several limitations remain.

\subsection{Limitations}

First, the experimental evaluation is conducted at a small to moderate scale, focusing on representative computational workflows rather than large, long-running scientific pipelines. While sufficient to validate execution semantics and reproducibility guarantees, these experiments do not capture the full complexity of production-scale scientific infrastructures.

Second, the current evaluation does not involve real laboratory hardware or cyber-physical systems. Although the framework is designed to support such settings through mediated action execution, interactions with physical instruments introduce additional safety, latency, and reliability considerations that are beyond the scope of this study.

Third, the framework relies on large language models as planning components, inheriting their limitations. Model behavior remains probabilistic, and the quality of proposed actions depends on the underlying LLM. While the execution layer constrains side effects and enforces determinism, the framework does not eliminate all sources of uncertainty arising from model-generated plans.

\subsection{Ethical Considerations}

This work does not claim autonomous scientific discovery or fully automated research. The framework is designed to assist, not replace, human scientific judgment. LAMs are treated as decision-support components whose outputs are mediated, constrained, and auditable, rather than as independent scientific actors.

The framework explicitly avoids unsafe experiment execution by requiring all actions to conform to predefined schemas and execution policies. Direct, unconstrained interaction with laboratory equipment or safety-critical systems is outside the intended scope of the system. Human oversight is required to define permissible actions, validate execution policies, and interpret results.

By prioritizing auditability, determinism, and provenance, the framework aims to reduce the risk of irreproducible or misleading scientific results arising from opaque automation. These design choices align with responsible AI principles for scientific research, emphasizing transparency, accountability, and human control over automated systems.
\section{Future Work}

While this work establishes a reproducibility-constrained execution framework for Large Action Models, several practical extensions remain for future investigation.

A natural next step is the integration of \emph{hardware-in-the-loop} execution. Extending the framework to interact with physical laboratory instruments or cyber-physical systems would require additional safety layers, latency handling, and fault isolation mechanisms. Such integration would allow the framework to be evaluated in real experimental environments while preserving reproducibility guarantees.

Another promising direction is \emph{multi-agent collaboration}. The current framework focuses on a single LAM interacting with a deterministic execution engine. Future work could explore coordinating multiple LAM instances operating over a shared or partially shared execution trace, enabling collaborative planning, parallel experimentation, or cross-validation of results. Maintaining consistent provenance and conflict resolution in such settings presents an open systems challenge.

Finally, future work may investigate \emph{formal verification of action traces}. While the framework enforces structured logging and deterministic execution, formal methods could be applied to verify properties of execution traces, such as completeness, causal consistency, or adherence to predefined safety constraints. Such verification could further strengthen trust in automated scientific workflows by providing machine-checkable guarantees over execution histories.

These directions aim to extend the framework’s applicability while preserving its core design principle: reproducibility as a first-class constraint in LAM-driven scientific automation.

\section{Conclusion}

Large Action Models offer a promising approach for automating complex scientific workflows by combining high-level reasoning with executable actions. However, unconstrained LAM-based systems are misaligned with the requirements of scientific research, where determinism, auditability, and reproducibility are fundamental. Without explicit safeguards, autonomous action execution risks producing results that cannot be independently verified or reliably reproduced.

In this paper, we introduced R-LAM, a reproducibility-constrained framework for deploying Large Action Models in scientific workflow automation. By formalizing actions, enforcing deterministic execution, and capturing complete provenance through execution trace graphs, the framework enables adaptive control while preserving reproducibility guarantees. Unlike prior approaches that emphasize either autonomous reasoning or static workflow execution, R-LAM bridges these paradigms by treating reproducibility as a core execution constraint.

We implemented R-LAM as a lightweight Python artifact and demonstrated its effectiveness through representative scientific workflows. The framework illustrates how LAM-driven automation can be integrated into scientific practice without compromising methodological rigor, supporting iterative experimentation, failure-aware execution, and auditable results.

We argue that reproducibility constraints are essential for deploying Large Action Models in scientific workflows.

\bibliographystyle{IEEEtran}
\bibliography{references}

\end{document}